\section{Move Semantics and Perfect Forwarding}\label{sec:move}
% Sources, in order: EMC++ Ch. 5, Items 23-30 -> Josuttis Ch. 4 (move semantics and container performance)
%   Most commonly asked modern-C++ topic.

\subsection{std::move and std::forward}
\src{EMC++ 23}
\begin{itemize}
  \item \texttt{std::move} is an unconditional cast to an rvalue ref, and \texttt{std::forward} only does this cast under specific conditions.
\end{itemize}
\begin{cppcode}
// std::move 
template<typename T>
typename remove_reference<T>::type&&
move(T&& param) {
  using ReturnType = typename remove_reference<T>::type&&;
  return static_cast<ReturnType>(param);
}
\end{cppcode}
\begin{itemize}
  \item \texttt{typename}: explicitly declares that a dependent qualified name is a type.
  \item \texttt{decltype(auto)}: deduce return type and keep refs
  \item the '\&\&' in the return type tells us that this reterns an rvalue ref.
  \item We must strip existing references since an rvalue ref to an lvalue ref collapses to an lvalue ref.
  \item Notice that there is no \texttt{remove\_cvref}; constness remains. 
  \item Careful!!! this can lead to silent copying.
\end{itemize}
\begin{cppcode}
struct Widget {
  Widget() = default;
  Widget(const Widget&) {std::cout << "copy ctor\n";}
  Widget(Widget&&) {std::cout << "move ctor\n";}
}
void test () {
  const Widget w;
  Widget w2 = std::move(w); // will invoke copy ctor
}
\end{cppcode}
\begin{itemize}
  \item A \texttt{const Widget\&\&} cannot bind to \texttt{Widget\&\&}, but can bind to \texttt{const Widget\&}.
  \item \cpp{std::forward} on the other hand is a conditional cast:
\end{itemize}
\begin{cppcode}
template <class T>
T&& forward(typename std::remove_reference_t<T>& t) noexcept;
\end{cppcode}
\begin{itemize}
  \item Recall that an rvalue ref to an lvalue ref collapses to an lvalue ref; forward utilizes this to ensure that function params preserve their original type.
  \item The act of passing a variable into a function will cause its value category to unconditinally change to an lvalue (since its given a name), even if its type is \cpp{T&&}.
  \item \cpp{std::forward} preserves the value category of that object before it was passed; if it was 
  \item concretely, "when t is a forwarding reference, this overload forwards the argument to another function with the value category it had when passed to the calling function".
  \item Forwarding reference: \cpp{T&&} is a forwarding reference iff the type T must be deduced when \cpp{T&&} is called.
\end{itemize}

\subsection{Universal References vs Rvalue References}
\src{EMC++ 24}
\begin{itemize}
  \item \cpp{T&&} has two meanings: one is an rvalue reference, and the other is either an rvalue or an lvalue (universal reference). 
\end{itemize}
\begin{cppcode}
void f(Widget&& param); // rvalue ref
Widget&& var1{Widget()}; // rvalue ref
auto&& var2 = var1; // not an rvalue ref (var1 is named, T& && collapses to T&)
template<typename T>
void f(std::vector<T>&& param); // rvalue ref (only binds to rvalue refs)
template<typename T>
void f(T&& param); // not rvalue ref (universal, can bind to either lvalue or rvalue))
\end{cppcode}
\begin{itemize}
  \item Universal references are permitted to bind to rvalue and lvalue references, and any cv qualifiers.
  \item These are also known as forwarding references as universal references should almost always be passed via \cpp{std::forward}.
  \item If an rvalue is passed into a function via a forwarding reference, it will implicitly be converted into an lvalue as it becomes named within the function body.
  \item This means that constructing an object here will always invoke the copy constructor; \cpp{std::forward<T>} ensures that the original value category that was passed is respected.
  \item passing an lvalue will remain unchanged, but passing an rvalue will invoke the move constructor.
  \item Universal references can also be a parameter pack; this is how \cpp{emplace_back} is implemented.
\end{itemize}
\begin{cppcode}
template<class T, class Allocator = allocator<T>>
class vector {
public:
template<class... Args>
void emplace_back(Args&&... args); // universal reference; when constructing the object in place it will do T(std::forward<Args>(args)...)
}
\end{cppcode}
\begin{itemize}
  \item Note that since universal references are when any deduction on the type has to be done in the call, \cpp{auto&&} is also a universal reference; this can be used in lambda functions to avoid copies.
\end{itemize}
\begin{cppcode}
auto timeFuncInvocation = 
  [](auto&& func, auto&&... params) { // universal references
    /* start timer... */
    std::forward<decltype(func)>(func)(
      std::forward<decltype<params>(params)...
    );
    /*stop timer and record time*/
  };
\end{cppcode}
\begin{itemize}
  \item Since we dont have access to \cpp{T} we need this \cpp{std::forward<decltype(arg)>} idiom:
  \item Has essentially the same behaviour as \cpp{std::forward<T>}.
\end{itemize}
consider:
\begin{cppcode}
  
\end{cppcode}

\subsection{std::move on Rvalue Refs, std::forward on Universal Refs}
\src{EMC++ 25}
\begin{itemize}
  \item \cpp{std::move} should only ever be used on rvalue references or if you are certain it is safe to steal resources.
  \item \cpp{std::forward} should be used whenever you have an object that may or may not be moved from. 
  \item I.e. rvalue refs should always be cast to rvalues (with move) when forwarding to other functions; universal references should sometimes be cast to rvalues (with forward) because theyre only sometimes rvalues.
\end{itemize}
\begin{cppcode}
class Widget {
public:
  Widget(Widget&& rhs) // rvalue ref; this can ONLY bind to rvalues
   : name(std::move(rhs.name)), p(std::move(rhs.p)) {}

  template<typename T>
  void setName(T&& newName) { //universal reference; can bind to anything
    name = std::forward<T>(newName); // conditional rvalue cast
  }
private:
  std::string name;
  std::shared_ptr<SomeDataStructure> p;

};
\end{cppcode}
\begin{itemize}
  \item Technically using \cpp{std::forward} on rvalue references will yield the same result as \cpp{std::move}:
\end{itemize}
\begin{cppcode}
void f(std::string&& s) {
  g(std::move(s)); // okay, as we have rvalue ref
  g(std::forward<std::string>(s)) // same result, just uglier
}
\end{cppcode}
\begin{itemize}
  \item Note that the other way around compiles but may lead to UB
\end{itemize}
\begin{cppcode}
class Widget {
public:
  template<typename T>
  void setName(T&& newName) { // universal; can bind to lvalue or rvalue
    name = std::move(newName); // bad!
  }
  /* ... */
};
Widget w;
std::string n = "my widget";
w.setName(n);
// using n is UB since it has been moved from
\end{cppcode}
\begin{itemize}
  \item While in the example above overloading \cpp{setName} with two variants (\cpp{const string&} and \cpp{string&&}) is an option, a function that takes $n$ arguments will require $2^n$ overloads.
  \item Functions that take an argument pack would require a seemingly infinite number of overloads; the universal reference ensures that only one implementation is required.
\end{itemize}
\begin{cppcode}
template<class T, class... Args>
shared_ptr<T> make_shared(Args&&... args);

template<class T, class... Args> 
unique_ptr<T> make_unique(Args&&... args);
\end{cppcode}
\begin{itemize}
  \item Since using after move is UB, forwarding a universal ref should be the final use of the reference (once its safe to steal its resources).
\end{itemize}
\begin{cppcode}
template<typename T>
void setSignText(T&& text) {
  sign.setText(text); // use text, but cannot steal yet
  auto now = std::chrono::system_clock::now();
  signHistory.add(now, std::forward<T>(text)); // conditionally cast to rvalue
}
\end{cppcode}
\begin{itemize}
  \item Returning an object passed as an rvalue or universal reference by value benefits from move or forwarding:
\end{itemize}
\begin{cppcode}
Matrix operator+(Matrix&& lhs, const Matrix& rhs) {
  lhs += rhs;
  return std::move(lhs); // doing return lhs will copy lhs into return value.
  
}
// similarly,
template<typename T>
requires std::convertible_to<T&&, Fraction> // C++20
Fraction reduceAndCopy(T&& frac) {
  frac.reduce();
  return std::forward<T>(frac);
}
\end{cppcode}
\begin{itemize}
  \item However, doing this on local variables returned by value is very bad!
\end{itemize}
\begin{cppcode}
Widget makeWidget() {
  Widget w;
  /* ... */
  w.doStuff();
  // return std::move(w); BAD!!
  return w;
}
\end{cppcode}
\begin{itemize}
  \item All modern compilers will perform an optimization called named return value optimization (NRVO) and construct \cpp{w} in the place of the return value slot.
  \item However, one of the requirements is that the returned local object must match exactly the return type in the function signature; since move casts to an rvalue ref, they dont match and NRVO cannot be performed.
\end{itemize}

\subsection{Avoid Overloading on Universal References}
\src{EMC++ 26}
Consider:
\begin{cppcode}
std::multiset<std::string> names;
void logAndAdd(const std::string& name) {
  auto now = std::chrono::system_clock::now();
  log(now, "logAndAdd");
  names.emplace(name);
}
/* ... */
std::string petName("Alice");
logAndAdd(petName); //lvalue string
logAndAdd(std::string("Bob")); //rvalue string
logAndAdd("AliceBob"); //string literal
\end{cppcode}
\begin{itemize}
  \item This however has inefficiencies when an rvalue is passed, since \cpp{name} is an lvalue it gets copied into \cpp{names}. 
  \item We can try using a univerasl reference to fix this:
\end{itemize}
\begin{cppcode}
template <typename T>
void logAndAdd(T&& name) {
  auto now = std::chrono::system_clock::now();
  log(now, "logAndAdd");
  names.emplace(std::forward<T>(name)); // moved for rvalues
}
\end{cppcode}
\begin{itemize}
  \item But what happens if we overload \cpp{logAndAdd}?
\end{itemize}
\begin{cppcode}
std::string nameFromIdx(int idx); // return name corresponding to idx

void logAndAdd(int idx)  {
  auto now = std::chrono::system_clock::now();
  log(now, "logAndAdd");
  names.emplace(nameFromIdx(idx));
}

short nameIdx{42};
logAndAdd(nameIdx); // ERROR
\end{cppcode}
\begin{itemize}
  \item Since the universal reference binds to everything (and is templated), it takes precedence over the int overload.
  \item We try to emplace into a string multiset and thus get an error.
  \item Universal reference functions are the greediest functions; they will take precedence over any overload that does not exactly match the signature.
  \item This can lead to some tricky side effects:
\end{itemize}
\begin{cppcode}
class Person {
public:
  template<typename T>
  Person(T&& n) : name(std::forward<T>(n)) {}
  
  Person(const Person& rhs); // copy ctor (compiler generated)
  Person(Person&& rhs); // move ctor (compiler generated)
};

Person p("Alice");
Person p2(p); //Error!
\end{cppcode}
\begin{itemize}
  \item Even though it looks like we are invoking the copy ctor, the type that p is when we pass it \cpp{Person&} matches the universal reference better (no const)
  \item Thus we call the constructor, and attempt to create a string from a Person type which errors.
\end{itemize}
\subsection{Alternatives to Overloading on Universal References}
\src{EMC++ 27}
\begin{itemize}
  \item Pre concepts: use tag dispatch
\end{itemize}
\begin{cppcode}
template<typename T>
void logAndAdd(T&& name) {
  logAndAddImpl(std::forward<T>(name), std::is_integral<std:remove_reference<T>::type());
}
\end{cppcode}
Note we have to add the \cpp{std::remove_reference} since if we pass an lvalue, T gets deduced to an lvalue reference and an int reference will return false for \cpp{std::integral}.

\begin{cppcode}
template<typename T>
void logAndAddImpl(T&&name, std::false_type) { // non-integral
  /* ... */
  names.emplace(std::forward<T>(name));
}

void logAndAddImpl(int idx, std::true_type) { logAndAdd(nameFromIdx(idx)); // guaranteed to match correct one}
\end{cppcode}
More modern, use \cpp{requires} or \cpp{if constexpr}
\begin{cppcode}
template<typename T>
  requires (!std::is_integral_V<std::remove_reference_t<T>>)
void logAndAdd(T&& name) {/* ... */}

// OR

// match everything to the universal ref
template <typename T>
void logAndAdd(T&& name) {
  if constexpr (std::is_integral_v<std::remove_reference_t<T>>) {
    /* ... */
  } else {
    /* ... */
  }
}
\end{cppcode}
The tricker case is when there is also inheritance:
\begin{cppcode}
class Person {
public:
    template <typename T>
    explicit Person(T&& n) : name(std::forward<T>(n)) {}
    Person(const Person& rhs);
    Person(Person&& rhs);
    ...
};

class SpecialPerson : public Person {
    SpecialPerson(const SpecialPerson& rhs) : Person(rhs) {}       // calls Person(T&&), not Person(const Person&)
    SpecialPerson(SpecialPerson&& rhs) : Person(std::move(rhs)) {} // same problem
};
\end{cppcode}
We can require that the forwarded ctor for Person is turned off if we pass a derived type and if it is not constructable from the type passed.
\begin{cppcode}
\begin{cppcode}
class Person {
public:
  template <typename T>
    requires (!std::derived_from<std::decay_t<T>, Person> && 
              std::constructible_from<std::string, T&&>)
    explicit Person(T&& n) : name(std::forward<T>(n)) {}
    ...
};
\end{cppcode}
the \cpp{std::decay_t<T>} strips any cv type quals and any references.
\subsection{Reference Collapsing}
\src{EMC++ 28}

\subsection{Assume Move Operations Are Not Present, Cheap, or Used}
\src{EMC++ 29}

\subsection{Perfect Forwarding Failure Cases}
\src{EMC++ 30}

\subsection{Moves in Containers}
\src{Josuttis Ch.~4}
